\begin{table}[t]
\centering
\caption{Representation and actor-free comparisons}
\label{tab:r5_resource}
\small
\begin{tabular}{rrrrr}
\toprule
$d$ & Actor-start loss & Zero-start loss & Quadratic loss range & Actor--zero difference \\
\midrule
4 & $6.446\times10^{-4}$ & $6.446\times10^{-4}$ & $2.065\times10^{-3}$--$3.203\times10^{-3}$ & $1.256\times10^{-10}$ \\
8 & $4.607\times10^{-4}$ & $4.607\times10^{-4}$ & $1.668\times10^{-3}$--$6.692\times10^{-3}$ & $1.104\times10^{-10}$ \\
16 & $5.665\times10^{-4}$ & $5.665\times10^{-4}$ & $5.444\times10^{-3}$--$6.777\times10^{-3}$ & $2.679\times10^{-11}$ \\
\bottomrule
\end{tabular}
\par\medskip
\begin{minipage}{.97\linewidth}
\footnotesize Each dimension includes seeds 101, 202, and 303 and the same 32 initial states per seed. The first two columns give the largest certified lifetime cost-loss upper bound across those seeds. The quadratic column spans the three maximum losses of a separately trained full positive-semidefinite quadratic critic. All use the same feasible controls and full shock trees; a reference Frank--Wolfe gap below $10^{-7}$ is added to each measured cost difference. Actor and zero starts use identical frozen nonlinear critic weights.
\end{minipage}
\end{table}
